\section{Main benchmark comparisons}
\label{sec:main-results-overview}
\begin{table}[!htbp]
\centering
\caption{Repair resolution rates (\%). Panel A tabulates
Figure~\ref{fig:benchmark-comparison}; the LoLBench comparison is in
Table~\ref{tab:new-codex55}. Panel B groups published reference results by
their reported configuration. Bold marks the best result per benchmark in
Panel A. An em dash indicates that the cited source does not report the
corresponding benchmark.}
\label{tab:main-results-overview}
\small
\setlength{\tabcolsep}{5pt}

\textbf{A. Three-benchmark comparison}\\[3pt]
\begin{tabular}{lccc}
\toprule
Method & Verified500 & Related-Lite99 & DeepSWE113 \\
\midrule
No memory & 61.6 & 21.4 & 60.8 \\
FAISS Flat & 65.2 & 26.8 & 64.6 \\
ExpeL & 62.4 & 20.7 & 58.7 \\
ExpeRepair$^\dagger$ & 63.6 & 31.3 & 64.4 \\
ReasoningBank & 67.0 & 24.9 & 61.7 \\
ACE$^\dagger$ & 62.2 & 28.1 & 62.8 \\
Supermemory & 63.6 & 26.5 & 66.2 \\
\textbf{ContextGraph} & \textbf{69.2} & \textbf{35.1} & \textbf{72.1} \\
\bottomrule
\end{tabular}

\smallskip
\begin{minipage}{\linewidth}\footnotesize
Verified500, Related-Lite99, and DeepSWE113 use DeepSeek V4 Pro 0813 and mini-SWE-agent.
Related-Lite99 and DeepSWE113 rates are means across five runs, each
reset to the same initial memory pool, with tasks evaluated in issue
creation-time order; per-run rates use completed evaluations as the
denominator.
ContextGraph adds lessons from completed successful and failed attempts
to memory for subsequent retrieval; baselines use their open-source
defaults for memory updates. The initial Verified500 memory draws on
Nebius SWE-bench-extra trajectories, which contain none of the evaluated
instances; the DeepSWE113 memory draws on trajectories from LoLBench
repositories (Appendix~\ref{sec:overlap-audit}).
$\dagger$ ExpeRepair and ACE are our adaptations to mini-SWE-agent.
Method sources: ExpeL~\citep{expel2024},
ExpeRepair~\citep{experepair2026}, ReasoningBank~\citep{reasoningbank2026},
ACE~\citep{ace2026}, and Supermemory~\citep{supermemoryreport2026}.
\end{minipage}

\medskip
\textbf{B. Published reference results, grouped by their reported configuration}\\[3pt]
\begin{tabular}{llrrl}
\toprule
Method & Backbone & \shortstack{Verified500\\(\%)} & \shortstack{Related-Lite99\\(\%)} & Source \\
\midrule
ExpeRepair & Claude 3.5 Sonnet & 57.2 & --- & ER \\
ExpeRepair & Claude 4 Sonnet & 74.6 & --- & ER \\
ReasoningBank & Gemini 2.5 Flash & 38.8 & --- & RB \\
ReasoningBank & Gemini 2.5 Pro & 57.4 & --- & RB \\
No memory & Claude Sonnet 4.5 & --- & 26.26 & SC \\
Supermemory & Claude Sonnet 4.5 & --- & 30.30 & SC \\
Oracle summary & Claude Sonnet 4.5 & --- & 34.34 & SC \\
\bottomrule
\end{tabular}

\smallskip
\begin{minipage}{\linewidth}\footnotesize
ER: \href{https://arxiv.org/html/2506.10484v2}{ExpeRepair, Table 1};
RB: \href{https://arxiv.org/html/2509.25140}{ReasoningBank, Table 2};
SC: \href{https://arxiv.org/html/2602.08316v3\#S3.T4}{SWE-ContextBench, Table 4}.
ER uses its own repair pipeline with occasional o4-mini fallback;
RB uses mini-SWE-agent; SC uses Claude Code.
Oracle summaries supply the known related history. Published scores use
their original budgets and memory settings;
Appendix~\ref{sec:published-swectx} reproduces the full published
SWE-ContextBench Lite table.
\end{minipage}
\end{table}
